% !TeX root = ../papers/nc-p3.tex
\section{第一長方形の反対階数と階数十の水平面の有限判定}
\label{sec:first-rectangle-rank-strata}
\begin{theorem}\label{thm:first-rectangle-rank-strata}
定理\ref{thm:first-rectangle-horizontal-frames}のトーラス開集合上で、反対階数は $8,10,14,15,16$ を取る。
$C=(u+v)(1+uv)$ と置き
\[
 F=u(v+1)^2z-C,\qquad G=v(u+1)^2z+C
\]
とする。$u=v=1$ で階数8、$u,v$ のちょうど一つが1なら階数10である。
両方とも1と異なるところでは階数は $16-\mathbf1_{F=0}-\mathbf1_{G=0}$ となる。
階数15の閉包は二つの有理曲面で、その共通部分の階数14の曲線は有理かつ滑らかである。
立ち入る定義域を上のトーラス開集合とすると、この曲線の閉包は階数10,8に達しない。
階数10の二面は階数8の直線で交わる。
\end{theorem}
\begin{proof}
反対切断の四色の四行列を取る。色0の行列式から単元を除くと $(v-1)^2F$、色1では $(u-1)^2G$、
色2,3では $(u-1)^2(v-1)^2$ である。三次の全小行列式のイデアルは四つの順に
$(v-1),(u-1),((u-1)(v-1)),((u-1)(v-1))$ で、各場合に一つの小行列式が生成元の単元倍となる。
全ての行列に可逆な二次小行列式がある。$u=1$ では $F=(v+1)^2(z-1)$ が非零であり、
$v=1$ では $G=(u+1)^2(z+1)$ が非零である。これで全ての階数を得る。

$F,G$ の各因子は $z$ について一次でそれぞれ有理グラフを与える。
$F=G=0$ なら $C+4uv=0$, $z=-4v/(v+1)^2$ である。
ここで $u,v\ne-1$ であり、$r=(u+1)/(v+1)$ と置けば
\[
 v=\frac{1-r}{r(r+1)},\qquad u=\frac{r(r-1)}{r+1}.
\]
逆写像も正則で、$r=0,\pm1$ はこの交叉にはない。従って曲線は有理直線の開部分である。
$u+1/u+v+1/v+4=0$ の勾配はこの部分では消えず、二グラフの交叉は横断的である。
$u=1$ または $v=1$ の先の非零条件が階数14の両面への特殊化を排除する。
一方 $F=0$ は $v=1$ の階数10面を $z=(u+1)^2/(4u)$ で切り、
$G=0$ は $u=1$ の面を $z=-(v+1)^2/(4v)$ で切る。
二面の交線 $u=v=1$ ではどちらも非零である。
\path{computations/verify_first_rectangle_rank_strata.py} が全小行列式と交線のパラメータを検算する。
\end{proof}

\begin{theorem}\label{thm:first-rectangle-edge-center}
階数10の両面では、折り畳み環は五つの中心四次式が生成する可換環上有限である。
従って各点について以下は標準Koszul左右AS正則性の必要十分条件である。
四次中心の基底を $s_1,\ldots,s_5$ とする。代数的に独立な $\lambda_1,\ldots,\lambda_4$ を加え、
四式 $s_j-\lambda_j s_5$ で割った各始点の長さ0から13の次元が $(1+T+T^2+T^3)^4$ の係数と一致することを検査する。
$u=-1$ または $v=-1$ の二面でも、符号を戻して同じ有限判定が全点に適用できる。
\end{theorem}
\begin{proof}
$u=1$ なら $t=b=1$, $c=v$, $x=(z+v)/(1+v)$ である。$v=1$ は既に正則と決定した交線であり、
階数10の面の基底環では $v z(v-1)(v+1)(v+z)(z-1)(vz+1)$ を逆にする。
全四次中心は文字0の三式と文字1の二式である。各位置の文字を $(0,2,2,0)$ だけ平行移動した表示で、
五式を全て割った商の経路を長さ9まで計算する。各始点で得られる開集合上の次元列は
\[
 (1,4,10,20,30,36,34,20,1,0)
\]
である。基底の決定と五次までの中心の持上げは単元だけで行え、以後の経路計算で必要になる追加の
既約因子の個数は始点順に $3,16,24,3$ となる。二位置移動は $(v,z)\mapsto(1/v,z)$ であり、
この像の四つの開集合も合わせる。

これら八集合のうち各点で少なくとも一つが使える。実際、始点0と3の因子から一つずつを選んで九つのイデアルを作る。
残る六リストの因子積を順に加え、上記の基底環で飽和すると九つとも単位イデアルになる。
ある始点 $p$ で五式の商の長さ9が零なら、他のどの始点からの経路も高々三歩で $p$ を通るので、
全商の長さ12は零である。従って元の環は五つの中心式上有限である。
\path{computations/verify_first_rectangle_edge_center.py} が八リストと九つの飽和を再計算する。
$v=1$ 側には一位置移動で移せる。

正則なら四次中心が生成する可換環の次元は4で、換算次数1の一般四式は斉次Noether正規化を与える。
定理\ref{thm:central-window-cutoff}の必要性から四式の商の次元列を得る。
逆に表示の列なら、中心商を持ち上げた自由加群の各次数が標準値である。
水平主成分の正則一般点からのHilbert下界と比較すると各始点で自由となり、
補題\ref{lem:parallel-residue-border}の最後の議論により全次数の左右AS正則性が従う。
負の二面は $t$ の符号変更の二コサイクル捻りであり、$u=v=-1$ の既知の交点も含めて判定できる。
\end{proof}

\begin{corollary}
76型の水平正則候補が少なくとも片側対称な表示を持つための必要十分条件は $t=1$ である。
階数10の二面ではそれぞれ一方の偶奇だけが対称化できる。
112型でも $u=1$ で偶位置交換、$v=1$ で奇位置交換を持つ表示が得られる。
\end{corollary}
\begin{proof}
$t=1$ は $u=1$ または $v=1$ と同値である。$b=1$ の補完は位置0,2の交換を保ち、
$c=1$ の補完は第四位置であらかじめ文字に1を加えれば位置1,3の交換を保つ。
これは両形の全係数での恒等式である。逆に片側対称なら反対切断は十次元の対称二乗を経由する。
76型の他の点の階数は14以上である。両側対称なら二位置移動で同型になるので、
周期4の階数10の二面では両側同時の対称化は起きない。
\end{proof}
